% ============================================================
% CHAPITRE 3 — SPRINT 0 : IDENTIFICATION DES BESOINS
% ============================================================

\chapter{Sprint 0 --- Identification des Besoins}

\section*{Plan}

\begin{tabular}{r l r}
\textbf{1} & \textbf{Identification des rôles} & \textbf{\pageref{sec:roles}} \\
\textbf{2} & \textbf{Identification des besoins} & \textbf{\pageref{sec:besoins}} \\
\textbf{3} & \textbf{Product Backlog} & \textbf{\pageref{sec:backlog}} \\
\textbf{4} & \textbf{Environnement de travail} & \textbf{\pageref{sec:environnement}} \\
\end{tabular}

\section*{Introduction}

Ce chapitre constitue la phase préparatoire du projet, correspondant au Sprint~0 dans le cadre méthodologique Scrum. Nous y définissons les rôles au sein de l'équipe Scrum, identifions de manière exhaustive les besoins fonctionnels et non fonctionnels du système, présentons le diagramme de cas d'utilisation global, élaborons le Product Backlog, et décrivons l'environnement de travail retenu --- tant sur le plan logiciel que matériel.

% ────────────────────────────────────────────────────────────
\section{Identification des rôles}
\label{sec:roles}
% ────────────────────────────────────────────────────────────

L'équipe Scrum constituée pour ce projet est composée des membres suivants~:

\begin{figure}[H]
    \centering
    \fbox{\parbox{10cm}{\centering\vspace{2cm}\textit{Diagramme de l'équipe Scrum}\vspace{2cm}}}
    \caption{Composition de l'équipe Scrum}
    \label{fig:scrum_team}
\end{figure}

\begin{table}[H]
\centering
\caption{Rôles de l'équipe Scrum}
\label{tab:scrum_roles}
\begin{tabularx}{\textwidth}{|l|l|X|}
\hline
\textbf{Rôle Scrum} & \textbf{Membre} & \textbf{Responsabilités} \\
\hline
Product Owner & TIM Tunisie & Définition des priorités, validation des incréments \\
\hline
Scrum Master & Mme Salsabil Hadji & Supervision académique, levée des obstacles \\
\hline
Développeur & Houcine Laghmani & Conception, développement, tests de tous les modules \\
\hline
\end{tabularx}
\end{table}


% ────────────────────────────────────────────────────────────
\section{Identification des besoins}
\label{sec:besoins}
% ────────────────────────────────────────────────────────────

L'analyse du système a permis d'identifier \textbf{quatre acteurs principaux} et \textbf{cinquante besoins fonctionnels} organisés en douze modules, ainsi que \textbf{dix besoins non fonctionnels}.

\subsection{Identification des acteurs}

Le système LLM Dashboard interagit avec quatre catégories d'acteurs, correspondant aux rôles RBAC implémentés dans le module d'authentification~:

\begin{table}[H]
\centering
\caption{Acteurs du système LLM Dashboard}
\label{tab:acteurs}
\begin{tabularx}{\textwidth}{|l|c|X|}
\hline
\textbf{Acteur} & \textbf{Rôle RBAC} & \textbf{Description} \\
\hline
Super Administrateur & \texttt{super\_admin} & Administrateur global de la plateforme. Accès total à tous les tenants, gestion des utilisateurs, configuration système, gestion des contrats et de la facturation. \\
\hline
Administrateur Tenant & \texttt{tenant\_admin} & Responsable d'un tenant spécifique (ex~: QALITAS ou GMAO PRO). Gère les clés API, les règles de gouvernance et consulte les tableaux de bord de son périmètre. \\
\hline
Observateur & \texttt{viewer} & Utilisateur en lecture seule. Consulte les tableaux de bord, les rapports et les métriques de son tenant. \\
\hline
Agent Applicatif & \texttt{api\_key} & Application cliente qui interagit avec la plateforme via des clés API pour router ses appels LLM à travers le Gateway. \\
\hline
\end{tabularx}
\end{table}

\subsection{Besoins fonctionnels}

Les besoins fonctionnels sont organisés par module, reflétant l'architecture modulaire du système. La priorisation suit la méthode MoSCoW~: \textbf{Must} (indispensable), \textbf{Should} (important), \textbf{Could} (souhaitable).

\subsubsection{BF-1~: Authentification et gestion des accès (Module Auth)}

\begin{table}[H]
\centering
\caption{Besoins fonctionnels --- Module Authentification}
\label{tab:bf_auth}
\begin{tabularx}{\textwidth}{|c|X|c|}
\hline
\textbf{ID} & \textbf{Description} & \textbf{Priorité} \\
\hline
BF-1.1 & Le système doit permettre l'inscription et l'authentification des utilisateurs via JWT (access token + refresh token) & Must \\
\hline
BF-1.2 & Le système doit supporter un RBAC à 4 rôles avec isolation stricte par tenant & Must \\
\hline
BF-1.3 & Le système doit permettre la création, la rotation et la révocation de clés API & Must \\
\hline
BF-1.4 & Le système doit supporter les clés virtuelles avec budgets plafonds par clé & Should \\
\hline
BF-1.5 & Le système doit permettre la déconnexion avec blacklisting des tokens JWT dans Redis & Must \\
\hline
\end{tabularx}
\end{table}

\subsubsection{BF-2~: Gateway LLM et routage intelligent (Module Gateway)}

\begin{table}[H]
\centering
\caption{Besoins fonctionnels --- Module Gateway}
\label{tab:bf_gateway}
\begin{tabularx}{\textwidth}{|c|X|c|}
\hline
\textbf{ID} & \textbf{Description} & \textbf{Priorité} \\
\hline
BF-2.1 & Le système doit proxyer les appels LLM vers 5+ fournisseurs (OpenAI, Anthropic, Google, Groq, Cohere) & Must \\
\hline
BF-2.2 & Le système doit journaliser chaque appel (tokens, coût, latence, statut) dans TimescaleDB & Must \\
\hline
BF-2.3 & Le système doit évaluer les règles de gouvernance avant chaque appel & Must \\
\hline
BF-2.4 & Le système doit publier un événement \texttt{llm.call.completed} après chaque appel & Must \\
\hline
\end{tabularx}
\end{table}

\subsubsection{BF-3~: Gouvernance (Module Gateway --- Governance)}

\begin{table}[H]
\centering
\caption{Besoins fonctionnels --- Gouvernance}
\label{tab:bf_gov}
\begin{tabularx}{\textwidth}{|c|X|c|}
\hline
\textbf{ID} & \textbf{Description} & \textbf{Priorité} \\
\hline
BF-3.1 & Supporter 5 types de règles~: \texttt{rate\_limit}, \texttt{budget\_cap}, \texttt{model\_restriction}, \texttt{token\_limit}, \texttt{time\_restriction} & Must \\
\hline
BF-3.2 & Évaluer les règles en cascade et journaliser chaque décision & Must \\
\hline
BF-3.3 & Permettre la gestion CRUD des règles par tenant & Must \\
\hline
BF-3.4 & Utiliser Redis comme compteur temps réel pour les quotas RPM & Should \\
\hline
\end{tabularx}
\end{table}

\subsubsection{BF-4~: Analytics et agrégation (Module Analytics)}

\begin{table}[H]
\centering
\caption{Besoins fonctionnels --- Module Analytics}
\label{tab:bf_analytics}
\begin{tabularx}{\textwidth}{|c|X|c|}
\hline
\textbf{ID} & \textbf{Description} & \textbf{Priorité} \\
\hline
BF-4.1 & Agréger les coûts quotidiennement par tenant, modèle et provider & Must \\
\hline
BF-4.2 & Calculer les rollups mensuels pour la facturation & Must \\
\hline
BF-4.3 & Supporter un pricing versionné avec plages de dates par modèle & Should \\
\hline
BF-4.4 & Exécuter les agrégations automatiquement via Celery Beat & Must \\
\hline
\end{tabularx}
\end{table}

\subsubsection{BF-5~: Détection d'anomalies (Module Detection)}

\begin{table}[H]
\centering
\caption{Besoins fonctionnels --- Module Détection}
\label{tab:bf_detection}
\begin{tabularx}{\textwidth}{|c|X|c|}
\hline
\textbf{ID} & \textbf{Description} & \textbf{Priorité} \\
\hline
BF-5.1 & Détecter les anomalies via 3 algorithmes ML~: STL, Isolation Forest, CUSUM & Must \\
\hline
BF-5.2 & Maintenir des baselines dynamiques par tenant/modèle/agent & Must \\
\hline
BF-5.3 & Classifier les anomalies par sévérité (info, warning, critical) & Should \\
\hline
BF-5.4 & Publier un événement \texttt{anomaly.detected} pour les traitements aval & Must \\
\hline
\end{tabularx}
\end{table}

\subsubsection{BF-6~: Explication IA (Module Explainer)}

\begin{table}[H]
\centering
\caption{Besoins fonctionnels --- Module Explainer}
\label{tab:bf_explainer}
\begin{tabularx}{\textwidth}{|c|X|c|}
\hline
\textbf{ID} & \textbf{Description} & \textbf{Priorité} \\
\hline
BF-6.1 & Générer une explication en langage naturel pour chaque anomalie & Should \\
\hline
BF-6.2 & Supporter plusieurs fournisseurs LLM pour la génération & Should \\
\hline
BF-6.3 & Produire 2--4 recommandations actionnables par anomalie & Should \\
\hline
BF-6.4 & Stocker les embeddings vectoriels pour la recherche sémantique (pgvector) & Could \\
\hline
\end{tabularx}
\end{table}

\subsubsection{BF-7~: Prévision budgétaire (Module Forecasting)}

\begin{table}[H]
\centering
\caption{Besoins fonctionnels --- Module Forecasting}
\label{tab:bf_forecast}
\begin{tabularx}{\textwidth}{|c|X|c|}
\hline
\textbf{ID} & \textbf{Description} & \textbf{Priorité} \\
\hline
BF-7.1 & Prévoir les tokens et coûts futurs selon 3 scénarios (optimiste, moyen, pessimiste) & Should \\
\hline
BF-7.2 & Identifier les risques de dépassement budgétaire et alerter & Should \\
\hline
\end{tabularx}
\end{table}

\subsubsection{BF-8~: Tableau de bord (Module Dashboard)}

\begin{table}[H]
\centering
\caption{Besoins fonctionnels --- Module Dashboard}
\label{tab:bf_dashboard}
\begin{tabularx}{\textwidth}{|c|X|c|}
\hline
\textbf{ID} & \textbf{Description} & \textbf{Priorité} \\
\hline
BF-8.1 & Fournir une vue exécutive avec KPIs (tokens, coût, requêtes, taux d'erreur) & Must \\
\hline
BF-8.2 & Afficher des graphiques de tendance par période (7j, 30j, 90j) & Must \\
\hline
BF-8.3 & Supporter les mises à jour en temps réel via WebSocket & Should \\
\hline
BF-8.4 & Intégrer un assistant IA conversationnel & Could \\
\hline
\end{tabularx}
\end{table}

\subsubsection{BF-9~: Facturation (Module Billing)}

\begin{table}[H]
\centering
\caption{Besoins fonctionnels --- Module Billing}
\label{tab:bf_billing}
\begin{tabularx}{\textwidth}{|c|X|c|}
\hline
\textbf{ID} & \textbf{Description} & \textbf{Priorité} \\
\hline
BF-9.1 & Supporter 3 types de contrats~: pay-as-you-go, forfait, hybride & Must \\
\hline
BF-9.2 & Générer automatiquement des factures mensuelles détaillées par modèle & Must \\
\hline
BF-9.3 & Supporter le cycle de vie des contrats~: brouillon $\rightarrow$ proposé $\rightarrow$ actif & Must \\
\hline
BF-9.4 & Produire des factures au format PDF professionnel & Should \\
\hline
\end{tabularx}
\end{table}

\subsubsection{BF-10~: Traçage (Module Tracing)}

\begin{table}[H]
\centering
\caption{Besoins fonctionnels --- Module Tracing}
\label{tab:bf_tracing}
\begin{tabularx}{\textwidth}{|c|X|c|}
\hline
\textbf{ID} & \textbf{Description} & \textbf{Priorité} \\
\hline
BF-10.1 & Tracer les sessions de conversation LLM avec corrélation parent-enfant & Should \\
\hline
BF-10.2 & Supporter l'A/B testing de prompts & Could \\
\hline
BF-10.3 & Permettre la configuration d'alertes multi-canal & Should \\
\hline
\end{tabularx}
\end{table}

\subsubsection{BF-11~: Gestion des prompts (Module Prompts)}

\begin{table}[H]
\centering
\caption{Besoins fonctionnels --- Module Prompts}
\label{tab:bf_prompts}
\begin{tabularx}{\textwidth}{|c|X|c|}
\hline
\textbf{ID} & \textbf{Description} & \textbf{Priorité} \\
\hline
BF-11.1 & Fournir un CMS de templates de prompts avec versionnement & Should \\
\hline
BF-11.2 & Supporter les variables Jinja2 dans les templates & Should \\
\hline
BF-11.3 & Gérer le cycle de publication~: brouillon $\rightarrow$ publié $\rightarrow$ archivé & Could \\
\hline
\end{tabularx}
\end{table}


\subsection{Besoins non fonctionnels}

\begin{table}[H]
\centering
\caption{Besoins non fonctionnels}
\label{tab:bnf}
\begin{tabularx}{\textwidth}{|c|l|X|}
\hline
\textbf{ID} & \textbf{Catégorie} & \textbf{Description} \\
\hline
BNF-1 & Performance & Latence du proxy Gateway $<$~50\,ms. Dashboard $<$~2 secondes. \\
\hline
BNF-2 & Scalabilité & Ajout de tenants et de modèles LLM sans redéploiement. \\
\hline
BNF-3 & Sécurité & Chiffrement PBKDF2-SHA256, hachage BLAKE2b, TLS, protection XSS. \\
\hline
BNF-4 & Disponibilité & Health checks HTTP, reconnexion automatique, arrêt gracieux. \\
\hline
BNF-5 & Multi-tenant & Isolation stricte des données par \texttt{tenant\_id}. \\
\hline
BNF-6 & Observabilité & Métriques Prometheus, logging JSON structuré. \\
\hline
BNF-7 & Maintenabilité & Clean Architecture, principes SOLID, couverture de tests. \\
\hline
BNF-8 & Portabilité & Conteneurisation Docker complète. \\
\hline
BNF-9 & Auditabilité & Journal d'audit append-only. \\
\hline
BNF-10 & i18n & Interface français / anglais. \\
\hline
\end{tabularx}
\end{table}


\subsection{Diagramme de cas d'utilisation global}

Le diagramme de cas d'utilisation global ci-dessous illustre les interactions entre les quatre acteurs identifiés et les principaux cas d'utilisation du système.

\begin{figure}[H]
    \centering
    \fbox{\parbox{14cm}{\centering\vspace{6cm}\textit{Diagramme de cas d'utilisation global de LLM Dashboard}\\[0.5cm]\textit{(4 acteurs~: Super Admin, Tenant Admin, Viewer, Agent API)}\\[0.5cm]\textit{(Cas d'utilisation~: S'authentifier, Consulter le dashboard, Gérer la gouvernance,}\\
    \textit{Consulter les anomalies, Gérer la facturation, Envoyer un appel LLM, etc.)}\vspace{3cm}}}
    \caption{Diagramme de cas d'utilisation global de LLM Dashboard}
    \label{fig:uc_global}
\end{figure}


% ────────────────────────────────────────────────────────────
\section{Product Backlog}
\label{sec:backlog}
% ────────────────────────────────────────────────────────────

Le Product Backlog ci-dessous recense l'ensemble des User Stories du projet, organisées par Epic et affectées aux sprints correspondants~:

\begin{longtable}{|c|c|p{7.5cm}|c|}
\hline
\textbf{Epic} & \textbf{ID} & \textbf{User Story} & \textbf{Sprint} \\
\hline
\endfirsthead
\hline
\textbf{Epic} & \textbf{ID} & \textbf{User Story} & \textbf{Sprint} \\
\hline
\endhead
\hline
\endfoot

Auth & US-1 & En tant qu'utilisateur, je veux m'authentifier pour accéder à la plateforme & 1 \\
\hline
Auth & US-2 & En tant qu'admin, je veux gérer les clés API pour contrôler l'accès des applications & 1 \\
\hline
Gateway & US-3 & En tant qu'agent, je veux envoyer un appel LLM via le proxy pour bénéficier du monitoring & 1 \\
\hline
Gateway & US-4 & En tant qu'admin, je veux configurer les règles de gouvernance pour contrôler les coûts & 1 \\
\hline
Dashboard & US-5 & En tant que manager, je veux consulter le tableau de bord exécutif pour suivre les KPIs & 1 \\
\hline
Dashboard & US-6 & En tant que viewer, je veux voir les graphiques de tendance pour analyser l'évolution & 1 \\
\hline
Detection & US-7 & En tant qu'admin, je veux être alerté des anomalies de consommation pour réagir rapidement & 2 \\
\hline
Explainer & US-8 & En tant que manager, je veux comprendre les causes d'une anomalie grâce à l'IA & 2 \\
\hline
Forecasting & US-9 & En tant que décideur, je veux prévoir les coûts futurs pour planifier mon budget & 2 \\
\hline
Analytics & US-10 & En tant qu'admin, je veux consulter les coûts agrégés par modèle et par jour & 2 \\
\hline
Optimizer & US-11 & En tant qu'admin, je veux recevoir des recommandations d'optimisation des coûts & 2 \\
\hline
Billing & US-12 & En tant que super admin, je veux gérer les contrats de facturation des tenants & 3 \\
\hline
Billing & US-13 & En tant que super admin, je veux générer automatiquement les factures mensuelles & 3 \\
\hline
Tracing & US-14 & En tant qu'admin, je veux tracer les sessions de conversation LLM & 3 \\
\hline
Prompts & US-15 & En tant qu'admin, je veux gérer un catalogue de templates de prompts & 3 \\
\hline
Tenants & US-16 & En tant que super admin, je veux gérer les tenants de la plateforme & 3 \\
\hline
Deploy & US-17 & En tant que DevOps, je veux déployer la plateforme via Docker Compose & 3 \\
\hline
\end{longtable}


% ────────────────────────────────────────────────────────────
\section{Environnement de travail}
\label{sec:environnement}
% ────────────────────────────────────────────────────────────

\subsection{Architecture logicielle}

La plateforme LLM Dashboard adopte une \textbf{architecture trois-tiers} (\textit{Three-Tier Architecture})~\cite{martin2018}~:

\begin{enumerate}[noitemsep]
    \item \textbf{Couche de présentation}~: application web Next.js~16 (React~19) avec TypeScript~;
    \item \textbf{Couche de logique métier}~: API REST FastAPI (Python~3.13) avec architecture modulaire~;
    \item \textbf{Couche de données}~: PostgreSQL~15 avec TimescaleDB + Redis~7.
\end{enumerate}

\begin{table}[H]
\centering
\caption{Stack technologique complète}
\label{tab:stack}
\begin{tabularx}{\textwidth}{|p{2.8cm}|p{3.5cm}|X|}
\hline
\textbf{Couche} & \textbf{Technologie} & \textbf{Justification} \\
\hline
Frontend & Next.js~16, React~19 & SSR/CSR hybride, routing natif, écosystème React \\
\hline
State Mgmt & React Query & Cache côté client, synchronisation, retry auto \\
\hline
Visualisation & Recharts~3.8 & Graphiques React performants et déclaratifs \\
\hline
Backend API & FastAPI (Python~3.13) & Performance async, validation Pydantic, Swagger auto~\cite{fastapi2024} \\
\hline
ORM & SQLAlchemy~2.x (async) & ORM mature, support async, migrations Alembic \\
\hline
Base de données & PostgreSQL~15 & Fiabilité ACID, types riches (JSONB, UUID)\\
\hline
Séries temporelles & TimescaleDB & Hypertables, compression, agrégations rapides~\cite{timescaledb2024}  \\
\hline
Vecteurs & pgvector & Embeddings pour la recherche sémantique \\
\hline
Cache / Broker & Redis~7 & Cache mémoire, pub/sub, compteurs atomiques \\
\hline
Tâches async & Celery~5.x & Orchestration de tâches avec scheduling~\cite{celery2024} \\
\hline
Reverse Proxy & Nginx & TLS, rate limiting, WebSocket, compression \\
\hline
Conteneurisation & Docker & Isolation, reproductibilité, déploiement \\
\hline
CI/CD & GitHub Actions & Pipeline automatisé, intégration GHCR \\
\hline
\end{tabularx}
\end{table}

\subsection{Architecture matérielle (infrastructure de déploiement)}

L'infrastructure de production est entièrement conteneurisée via Docker Compose et comprend \textbf{six services}~:

\begin{table}[H]
\centering
\caption{Services Docker de la plateforme}
\label{tab:docker_services}
\begin{tabularx}{\textwidth}{|l|l|X|}
\hline
\textbf{Service} & \textbf{Image} & \textbf{Rôle} \\
\hline
\texttt{postgres} & \texttt{timescale/timescaledb:latest-pg15} & Base de données avec extensions TimescaleDB et pgvector \\
\hline
\texttt{redis} & \texttt{redis:7-alpine} & Cache, compteurs de quotas, broker Celery \\
\hline
\texttt{api} & Image custom (Dockerfile) & Serveur API FastAPI (Uvicorn, 4 workers) \\
\hline
\texttt{worker} & Image custom (Dockerfile) & Worker Celery pour les tâches asynchrones \\
\hline
\texttt{scheduler} & Image custom (Dockerfile) & Celery Beat pour les tâches périodiques \\
\hline
\texttt{nginx} & \texttt{nginx:alpine} & Reverse proxy, terminaison TLS, rate limiting \\
\hline
\end{tabularx}
\end{table}

\begin{figure}[H]
    \centering
    \fbox{\parbox{14cm}{\centering\vspace{4cm}\textit{Diagramme de déploiement Docker}\\[0.5cm]\textit{(6 services : PostgreSQL + Redis + API + Worker + Scheduler + Nginx)}\\[0.5cm]\textit{(Réseau interne : llm-net / Volumes persistants : postgres\_data, redis\_data)}\vspace{3cm}}}
    \caption{Architecture de déploiement Docker Compose}
    \label{fig:docker_arch}
\end{figure}

\subsection{Modèle de données physique}

Le schéma de base de données de la plateforme comprend \textbf{19 tables} organisées en domaines fonctionnels. Le diagramme entité-relation ci-dessous présente les entités principales et leurs relations~:

\begin{figure}[H]
    \centering
    \fbox{\parbox{14cm}{\centering\vspace{5cm}\textit{Diagramme Entité-Relation --- Modèle physique}\\[0.5cm]\textit{(19 tables : llm\_tenant, llm\_auth\_user, llm\_api\_key, llm\_token\_log,}\\
    \textit{governance\_rule, llm\_cost\_daily, llm\_cost\_monthly, llm\_model\_pricing,}\\
    \textit{llm\_anomaly, anomaly\_explanation, llm\_token\_baseline, forecast\_result,}\\
    \textit{billing\_contract, invoice, invoice\_line\_item, llm\_session, prompt\_template, ...)}\\[0.5cm]\textit{(Voir fichier PlantUML~: ch3\_er\_database.puml)}\vspace{3cm}}}
    \caption{Modèle de données physique de la plateforme LLM Dashboard}
    \label{fig:er_database}
\end{figure}



\section*{Conclusion}

Ce chapitre a permis de définir le cadre complet du Sprint~0. L'identification des rôles Scrum, la spécification exhaustive de 50 besoins fonctionnels répartis en 11 modules, la définition de 10 besoins non fonctionnels, la modélisation du cas d'utilisation global et l'élaboration du Product Backlog fournissent une base solide pour les sprints de développement. L'environnement de travail, fondé sur une stack technologique moderne et une infrastructure Docker conteneurisée, garantit la reproductibilité et la scalabilité de la solution. Le chapitre suivant détaille le Sprint~1, consacré à la mise en place de l'infrastructure technique et des modules fondamentaux.
